% ============================================================
%  深圳大学实验报告 —— 项目一：图像编码器前向计算复现
%  版式参照课程提供的实验报告模板；编译方式：Overleaf + XeLaTeX
%  结构：封面 1 页 + 正文 2 页 + 批阅页 1 页，共 4 页
% ============================================================

\documentclass[a4paper,12pt]{article}

% ----------------------------
% 中文与字体支持
% ----------------------------
\usepackage{fontspec}
\usepackage{xeCJK}
\setCJKmainfont{Noto Serif CJK SC}

% ----------------------------
% 页面设置与常用宏包
% ----------------------------
\usepackage{geometry}
\geometry{left=0.85in, right=0.85in, top=0.8in, bottom=0.8in}

\usepackage{graphicx}
\usepackage{fancyhdr}
\usepackage{tikz}
\usetikzlibrary{calc, positioning, arrows.meta}
\usepackage{float}
\usepackage{amsmath}
\usepackage{array}
\usepackage{booktabs}
\usepackage{enumitem}
\usepackage{listings}
\usepackage{xcolor}
\usepackage[hidelinks]{hyperref}

\setlist{nosep, leftmargin=1.6em, itemsep=1pt, parsep=0pt}
\setlength{\parskip}{2pt}
\setlength{\floatsep}{5pt}
\setlength{\textfloatsep}{5pt}
\setlength{\intextsep}{5pt}

\fancyhf{}

% ----------------------------
% 代码高亮样式
% ----------------------------
\definecolor{codebg}{RGB}{248,248,251}
\definecolor{codekw}{RGB}{0,82,150}
\definecolor{codecom}{RGB}{0,120,80}
\definecolor{codestr}{RGB}{160,55,20}
\definecolor{gray65}{RGB}{105,105,105}

\lstdefinestyle{pystyle}{
    language=Python,
    backgroundcolor=\color{codebg},
    basicstyle=\ttfamily\footnotesize,
    keywordstyle=\color{codekw}\bfseries,
    commentstyle=\color{codecom}\itshape,
    stringstyle=\color{codestr},
    frame=lines, framesep=2pt, rulecolor=\color{gray65},
    breaklines=true, breakatwhitespace=false,
    showstringspaces=false,
    columns=fullflexible, keepspaces=true,
    aboveskip=3pt, belowskip=3pt,
}
\lstset{style=pystyle}

\newcommand{\code}[1]{{\ttfamily\footnotesize #1}}

% ============================================================
\begin{document}
\linespread{0.95}\selectfont   % 略微收紧行距，正文控制在 2 页内

% ============================================================
% 封面页
% ============================================================
\begin{titlepage}
    \centering
    \vspace*{2cm}
    \Huge{\textbf{深 \ 圳 \ 大 \ 学 \ 实 \ 验 \ 报 \ 告}}\\[1.5cm]

    {\Large\textbf{课程名称}\quad\underline{\makebox[8cm][c]{Python程序设计}}}\\[0.5cm]
{\Large
\noindent\makebox[3.2cm][l]{\textbf{项目名称}}
\underline{\makebox[8cm][c]{项目一 \ 图像编码器前向计算复现}}}\\[0.5cm]
{\Large
\noindent\makebox[3.2cm][l]{\textbf{学 \quad \quad 院}}
\underline{\makebox[8cm][c]{人工智能学院}}}\\[0.5cm]
{\Large
\noindent\makebox[3.2cm][l]{\textbf{专 \quad \quad 业}}
\underline{\makebox[8cm][c]{人工智能}}}\\[0.5cm]
{\Large
\noindent\makebox[3.2cm][l]{\textbf{指导教师}}
\underline{\makebox[8cm][c]{樊超}}}\\[0.5cm]
{\Large
\noindent\makebox[3.2cm][l]{\textbf{报 \quad 告 \quad 人}}
\underline{\makebox[3cm][c]{贾潘}}
\hspace{0.5cm}
\textbf{学号}
\underline{\makebox[3cm][c]{2025681051}}}\\[0.5cm]
{\Large
\noindent\makebox[3.2cm][l]{\textbf{实验时间}}
\underline{\makebox[8cm][c]{2026/11/19}}}\\[0.5cm]
{\Large
\noindent\makebox[3.2cm][l]{\textbf{提交时间}}
\underline{\makebox[8cm][c]{2026/11/19}}}\\[1.5cm]

    \vfill
    \Large{\textbf{教务处制}}
\end{titlepage}

\newpage

% ============================================================
% 正文
% ============================================================

\section{任务与实现}

本项目用 NumPy 复现 OpenCLIP ViT-B-32 的\textbf{图像编码器前向计算}：输入预处理后的
\code{[B,3,224,224]} 图像，输出每张图一个单位长度的 \code{[B,512]} 向量，并要求与独立
FP32 参考结果在给定容差内一致。按提交要求只编辑 \code{image\_encoder.py} 的六组 TODO，
其余框架、权重与数值辅助函数均不改动。六组实现与其 shape 契约如表~\ref{tab:todo}。

\begin{table}[H]
\centering
\footnotesize
\setlength{\tabcolsep}{3pt}
\caption{六组 TODO 的实现要点与 shape 变化}
\label{tab:todo}
\begin{tabular}{@{}l l p{9.2cm}@{}}
\toprule
\textbf{TODO} & \textbf{方法} & \textbf{核心运算与 shape 变化} \\
\midrule
01 & \code{patchify} & \code{[B,3,224,224]} $\to$ \code{[B,49,3072]}：拆成 $7\times7$ 个 $32\times32$ patch \\
02 & \code{split\_heads} & \code{[B,N,768]} $\to$ \code{[B,12,N,64]}：\code{reshape}+\code{transpose} \\
03 & attention & \code{[B,12,N,64]} $\to$ \code{[B,N,768]}：缩放点积、加权求和、合并头、输出投影 \\
04 & block & \code{[B,N,768]} $\to$ \code{[B,N,768]}：两次 LayerNorm、两次残差、GELU MLP \\
05 & \code{embed} & \code{[B,49,3072]} $\to$ \code{[B,50,768]}：无 bias 投影 + CLS + 位置编码 \\
06 & \code{\_\_call\_\_} & \code{[B,50,768]} $\to$ \code{[B,512]}：CLS、post-Norm、投影、单位化 \\
\bottomrule
\end{tabular}
\end{table}

\subsection{前向链路与一个 shape 变化}

图~\ref{fig:chain} 给出整条 shape 链路。最关键的一处 shape 变化是
\code{split\_heads}：注意力权重 \code{qkv} 是 \code{[B,N,2304]}，切成 q/k/v 后各为
\code{[B,N,768]}；要把它变成 12 个并行的头，必须先 \code{reshape(B,N,12,64)} 把特征轴
拆成「头数 $\times$ 头维度」，再 \code{transpose(0,2,1,3)} 把头轴移到第 1 维，得到
\code{[B,12,N,64]}。这样 \code{q @ k.transpose(0,1,3,2)} 才是合法的批量矩阵乘，得到
\code{[B,12,N,N]} 的分数；加权求和后按相反顺序 \code{transpose(0,2,1,3).reshape(B,N,768)}
合并回特征轴。

\begin{lstlisting}[caption={多头拆分与分数计算（\code{image\_encoder.py} 核心两行）}, label={lst:head}]
x.reshape(B, N, 12, 64).transpose(0, 2, 1, 3)          # [B,N,768] -> [B,12,N,64]
q @ k.transpose(0, 1, 3, 2) / np.sqrt(np.float32(64))   # -> [B,12,N,N]
\end{lstlisting}

\begin{figure}[H]
\centering
\begin{tikzpicture}[font=\footnotesize,
    blk/.style={draw=black!55, rounded corners=2pt, minimum width=2.55cm,
                minimum height=0.92cm, fill=blue!8, align=center},
    arr2/.style={-{Latex[length=1.6mm]}, thick, black!60}]
  \node[blk] (n1) at (0,0)     {图像\\ \scriptsize [B,3,224,224]};
  \node[blk] (n2) at (3.55,0)  {patch + embed\\ \scriptsize [B,50,768]};
  \node[blk, fill=green!12] (n3) at (7.10,0) {12 $\times$ Block\\ \scriptsize [B,50,768]};
  \node[blk, fill=violet!12] (n4) at (10.65,0) {CLS + 投影\\ \scriptsize [B,512]};
  \draw[arr2] (n1) -- (n2) node[midway, above=1mm, font=\scriptsize] {01, 05};
  \draw[arr2] (n2) -- (n3) node[midway, above=1mm, font=\scriptsize] {02--04};
  \draw[arr2] (n3) -- (n4) node[midway, above=1mm, font=\scriptsize] {06};
\end{tikzpicture}
\caption{图像编码器前向计算的 shape 链路，箭头上是对应的 TODO 编号}
\label{fig:chain}
\end{figure}

\section{结果与分析}

在 Windows 11 / Python 3.13.14 / NumPy 2.5.5 下运行 \code{python check.py -{}-stage all}，
八个阶段全部通过，逐阶段误差如表~\ref{tab:err}。输出向量容差 atol=$2\!\times\!10^{-5}$、
rtol=$2\!\times\!10^{-4}$，内部张量 atol=$2\!\times\!10^{-4}$；实测最大绝对误差
$4.387\times10^{-5}$（tokens），最终输出仅 $5.662\times10^{-7}$，均远小于容差。

\begin{table}[H]
\centering
\footnotesize
\setlength{\tabcolsep}{4pt}
\caption{\code{check.py -{}-stage all} 的逐阶段数值误差}
\label{tab:err}
\begin{tabular}{@{}l l r r l@{}}
\toprule
\textbf{阶段} & \textbf{shape} & \textbf{最大绝对误差} & \textbf{平均绝对误差} & \textbf{结果} \\
\midrule
preprocess   & (10,3,224,224) & 0.000e+00 & 0.000e+00 & PASS \\
patches      & (2,49,3072)    & 0.000e+00 & 0.000e+00 & PASS \\
tokens       & (2,50,768)     & 4.387e-05 & 1.335e-06 & PASS \\
block0       & (2,50,768)     & 7.147e-06 & 3.348e-07 & PASS \\
last\_tokens & (2,50,768)     & 4.196e-05 & 2.633e-06 & PASS \\
output       & (10,512)       & 5.662e-07 & 8.688e-08 & PASS \\
batch        & (3,512)        & 4.061e-07 & 7.848e-08 & PASS \\
\bottomrule
\end{tabular}
\end{table}

\subsection{一个修掉的 bug}

课堂示例 \code{transformer\_steps.py} 用二维表，转置写作 \code{k.T}；但真实
\code{k} 是四维 \code{[B,12,N,64]}，\code{.T} 会\textbf{反转全部轴}，导致
\code{q @ k.T} 报 “matmul: size 12 is different from 64”。改为
\code{k.transpose(0,1,3,2)} 只交换最后两个轴、保留 batch 与 head 组后，
\code{block0} 误差降到 $7.147\times10^{-6}$ 并通过。这也是把讲义算术迁移到真实四维
张量时最容易出错的一步。

\subsection{数值约定与讨论}

线性层权重按 \code{[out,in]} 存储，故 \code{qkv}、\code{out\_proj}、MLP 两处都要
\code{@ W.T}；但 \code{visual.proj} 已是 \code{[768,512]} 的 \code{[in,out]}，不能再转置。
patch 投影无 bias，scores 需除以 $\sqrt{64}=8$，激活与 LayerNorm 用随附的精确 GELU 与
$\varepsilon=10^{-5}$。全程保持 float32，参考输出仅用于校验、绝不作为编码器输入；
\code{batch} 阶段换序与重复图片同样通过，说明实现不依赖固定批量。

\section{提交信息}

\begin{itemize}
    \item \textbf{提交文件}：完成的 \code{image\_encoder.py}、简短 README、本报告；大权重文件与参考数据按要求不重复提交；
    \item \textbf{运行环境与参数}：Windows 11 / Python 3.13.14 / NumPy 2.5.5 / Pillow / imageio；ViT-B-32，patch $32\!\times\!32$，12 头 12 层，width 768，输出 512，全程 float32；
    \item \textbf{复现命令}：\code{python check.py -{}-stage all} 与 \code{python run.py -{}-output output\_tokens.npz}；
    \item \textbf{本人贡献与所用协助}：六组 TODO 的实现、逐阶段调试与数值核对由本人独立完成；实现依据为课程讲义、随附 README、\code{transformer\_steps.py} 示例及 \code{check.py} 校验工具。
\end{itemize}

\newpage

% ============================================================
% 批阅与成绩评定页
% ============================================================
\begin{tikzpicture}[remember picture, overlay]
  \draw[thick]
    ($(current page.north west)+(2cm,-3cm)$)
    rectangle
    ($(current page.south east)+(-2cm,3.5cm)$);
\end{tikzpicture}

\vspace{1cm}

\noindent \textbf{指导教师批阅意见：}
\vspace{6cm}
\hfill

\vspace{1cm}

\noindent \textbf{成绩评定：}
\vspace{2cm}
\hfill

\vspace{1cm}

\noindent
\hfill
\textbf{指导教师签字：}\quad\quad\quad\quad\quad\\[0.5cm]

\hfill
\textbf{年}\quad\quad
\textbf{月}\quad\quad
\textbf{日}\quad

\vspace{1cm}

\noindent \textbf{备注：}

\begin{tikzpicture}[remember picture, overlay]
  \node[
    anchor=north west,
    inner sep=0pt,
    text width=\dimexpr\paperwidth-4cm\relax
  ] at ($(current page.south west)+(2cm,3.1cm)$) {
    \begin{minipage}{\dimexpr\paperwidth-4cm\relax}
      \small
      \noindent
      \textbf{注：}\;1、报告内的项目或内容设置，可根据实际情况加以调整和补充。\\[0.2cm]
      \phantom{\textbf{注：}\;}2、教师批改学生实验报告时间应在学生提交实验报告时间后 10 日内。
    \end{minipage}
  };
\end{tikzpicture}

\end{document}
